\subsection{Classes}

Eg:
\begin{mdframed}[nobreak=true,backgroundcolor=gray!10,linecolor=Firebrick4]
\begin{verbatim}
class Point {
public:
    Point(int x, int y) : x_(x), y_(y) {}
private:
    int x_;
    int y_;
};
\end{verbatim}
\end{mdframed}

\subsubsection{Inheritance Behaviour}

\begin{mdframed}[nobreak=true,backgroundcolor=SlateGray2!40,linecolor=Firebrick4]
\textbf{Simple rule}
\begin{enumerate}
\item Private members are never inherited.
\item \texttt{derived-type = min(base-type, inheritance-type)}.
\end{enumerate}
\end{mdframed}

\begin{tabular}{|l|l|l|l|}
\hline
Member / Inheritance & Public    & Protected & Private \\ \hline
Public               & Public    & Protected & Private \\ \hline
Protected            & Protected & Protected & Private \\ \hline
Private              & NA        & NA        & NA      \\ \hline 
\end{tabular}

\subsubsection{Members Variables}

\begin{itemize}
\item \texttt{this} - pointer to the current object.
\item \texttt{static} - shared across all objects of the class.
\end{itemize}




\subsubsection{Member Functions}

\begin{center}
\textbf{Rule-of-5 members} 
\end{center}

\textbf{1. Copy Ctor}
\begin{mdframed}[nobreak=true,backgroundcolor=gray!10,linecolor=Firebrick4]
\begin{verbatim}
// Pass by ref to avoid recursion.
Vec::Vec(const Vec &other) { 
    n_ = other.n_;
    data_ = new int[n_];
    for (int i = 0; i < n_; i++) {
        data_[i] = other.data_[i];
    }
}    
\end{verbatim}
\end{mdframed}

\textbf{2. Move Ctor}

\begin{mdframed}[nobreak=true,backgroundcolor=gray!10,linecolor=Firebrick4]
\begin{verbatim}
Vec::Vec(Vec &&other) noexcept {
    n_ = other.n_;
    data_ = other.data_;
    other.n_ = 0;
    other.data_ = nullptr;    
}    
\end{verbatim}
\end{mdframed}

\vfill\null
\columnbreak


\textbf{3. Copy Assignment Operator}

\begin{mdframed}[nobreak=true,backgroundcolor=gray!10,linecolor=Firebrick4]
\begin{verbatim}
Vec& Vec::operator=(const Vec &other) {
    if (this != &other) { // Self-assignment check.
        delete[] data_; // Free existing resource.
        n_ = other.n_;
        data_ = new int[n_];
        // Deep copy.
        for (int i = 0; i < n_; i++) {
            data_[i] = other.data_[i];
        }
    }    
    return *this; // To permit chained assignments.
}
\end{verbatim}
\end{mdframed}


\textbf{4. Move Assignment Operator}

\begin{mdframed}[nobreak=true,backgroundcolor=gray!10,linecolor=Firebrick4]
\begin{verbatim}
Vec& Vec::operator=(Vec &&other) noexcept {
    if (this != &other) { // Self-assignment check.
        delete[] data_; // Free existing resource.
        n_ = other.n_;
        data_ = other.data_;
        other.n_ = 0;
        other.data_ = nullptr;    
    }    
    return *this; // To permit chained assignments.    
}
\end{verbatim}
\end{mdframed}

\textbf{5. Destructor}

\begin{mdframed}[nobreak=true,backgroundcolor=gray!10,linecolor=Firebrick4]
\begin{verbatim}
Vec::~Vec() {
    delete[] data_;
}    
\end{verbatim}
\end{mdframed}

\begin{center}
\textbf{Others}
\end{center}

\textbf{6. Default Ctor}
Eg: \texttt{Foo::Foo(){}}

\textbf{7. Converting Ctor}\\
Ctor with single argument or multiple arguments with default values.\\
\begin{mdframed}[nobreak=true,backgroundcolor=gray!10,linecolor=Firebrick4]
Say: \texttt{Foo::Foo(const Bar\&);}
\begin{verbatim}
Foo f = bar; // Implicit conversion from Bar to Foo.
\end{verbatim}
\texttt{explicit Foo(const Bar\&);// Prevents implicit conversion.} 
\end{mdframed}

\textbf{8. Static methods}\\
Belong to class, not object.\\
No \texttt{this} pointer.\\


\vfill\null
\columnbreak


\begin{mdframed}[nobreak=true,backgroundcolor=magenta!10,linecolor=magenta]
Compilers provided default impelementation for 1-6 methods.
\begin{itemize}
\item \textbf{\texttt{delete}}: Prevents default generation.
\item \textbf{\texttt{default}}: Makes the intent clear.\\
                                 \textcolor{red}{Forces default generation?}
\item \textbf{\texttt{explicit}}: Prevents implicit conversions for converting ctors.
\end{itemize}
\end{mdframed}


\subsubsection{Virtual Methods and Abstract Class}

\begin{itemize}
\item \textbf{Virtual method: } Eg: \texttt{virtual void Animal::speak();}
\item \textbf{Pure virtual method: } Eg: \texttt{virtual void Animal::speak() = 0;}
\item \textbf{Abstract class:} Class with at least one pure virtual method.\\
                               Cannot instantiate objects of abstract class.
\end{itemize}

\begin{mdframed}[nobreak=true,backgroundcolor=magenta!10,linecolor=magenta]
\textbf{Virtual vs normal methods:}\\
\begin{verbatim}
Animal* a = new Dog();
a->speak();
\end{verbatim}
\begin{itemize}
\item \texttt{speak} is virtual: Dog's speak() called.
\item \texttt{speak} is normal: Animal's speak() called.
\end{itemize}
\end{mdframed}


\subsubsection{override \& final}

Specifiers at the end. Relative order between them doesn't matter.\\

\textbf{\texttt{override}}:\\
Tells a method is meant to override a base class method. \\
Compiler error if it doesn't.\\
Eg:
\begin{mdframed}[nobreak=true,backgroundcolor=gray!10,linecolor=Firebrick4]
\textbf{Base:} \texttt{void Animal::speak() const;}\\
\textbf{Derived:} \texttt{void Dog::speak(); // Compiles.}\\
\textbf{Unexpected runtime behaviour:} \\
\texttt{a->speak(); // Base class speak.}\\
\textbf{Derived w/ override:} \\
\texttt{void Dog::speak() override; // Compiler error.}\\
\end{mdframed}

\textbf{\texttt{final}}:\\
Method not supposed to be overridden in derived classes.\\
Compiler error if it is.\\

Eg:
\begin{mdframed}[nobreak=true,backgroundcolor=gray!10,linecolor=Firebrick4]
\textbf{Base:} \texttt{void Dog::speak() override final;}\\
\textbf{Derived:} \\
\texttt{void BorderCollie::speak() override; // Compiler error.}
\end{mdframed}